% One-time imported draft. Editable native source; no Markdown regeneration.
\section{Step 16 Explain the interpretation and width
  rules}\label{004-bit-inspector-s16-step-16-explain-the-interpretation-and-width-rules}

The program separates a stored word, two ways to extend it, two
arithmetic widths, and four ways to display the original word. Keeping
those contracts distinct makes boundary behavior predictable.

\subsection{Write or
  inspect}\label{004-bit-inspector-s16-write-or-inspect}

Run the public checks and answer the practice questions. Include an
eight-byte result trace and explain why poisoning the output fields
exposes a narrow store.

\subsection{Expected
  observation}\label{004-bit-inspector-s16-expected-observation}

The final grader accepts different registers and instruction sequences,
but every output byte and all report text must match for substituted
input words.

\subsection{Check your
  understanding}\label{004-bit-inspector-s16-check-your-understanding}

What determines whether the high register half must be checked?

\begin{itemize}
  \tightlist
  \item
        \textbf{A} Whether the default input happens to be positive.
  \item
        \textbf{B} Whether the source uses short register names.
  \item
        \textbf{C} The full-width result contract, even when printed output
        uses only a word.
\end{itemize}
